% GENERATED FILE. Edit canonical lessons, then run npm run generate:books.
% canonical-source-hash: fnv1a64:47d7539e2b9461e7
% canonical-lessons: FR-C470-frayeur, FR-C470-effrayer, FR-C470-sursauter, FR-C470-soulagement, FR-C470-rassurer

\chapter{Quelle frayeur !}
\label{ch:fr-quelle-frayeur}
\hlchaptermodality{470}

\begin{chapteropening}
\textbf{By the end of this chapter:} \emph{I can talk about a fright and the relief afterwards.}

Complete the last lesson of chapter 470: I can talk about a fright and the relief afterwards.
\end{chapteropening}

\section[la frayeur]{la frayeur — fright, scare}

\label{lesson:FR-C470-frayeur}



\begin{quote}
Before the new one: say the French for to annoy, then the French for to irritate.
\end{quote}

\subsection*{You'll want to know: la frayeur}
\textbf{la frayeur} — \textquotedblleft{}fright, scare\textquotedblright{}.

\begin{quote}
Quelle frayeur ! J'ai cru que j'avais perdu mon passeport.
\end{quote}

\begin{quote}
\emph{What a fright! I thought I had lost my passport.}
\end{quote}

Stronger than \emph{la peur}, and over quickly.

\begin{cousinweb}
\emph{la frayeur} is inherited from Old French \emph{freor} (EtymDB), from Latin \emph{fragor} \textquotedblleft{}a crash, a din\textquotedblright{}, built on \emph{frangere} \textquotedblleft{}to break\textquotedblright{} (Lewis \& Short); its form was later drawn towards \emph{effrayer}.
\end{cousinweb}

\begin{grammarlens}[title={What you've built}]
One more word for telling, explaining and describing at B1.
\end{grammarlens}

\subsection*{Your turn}
\begin{itemize}
\raggedright
  \item \emph{Say it:} \emph{la frayeur}
  \item \emph{Say it:} \emph{la frayeur}, once more
  \item \emph{Say it:} a sentence of your own with \emph{la frayeur}
  \item \emph{Recall:} say the French for to annoy, then the French for to irritate, then say \emph{la frayeur} again
\end{itemize}

\begin{tcolorbox}[breakable,colback=teal!4,colframe=teal!35!black,title={Before you move on}]
What does la frayeur mean? (\textquotedblleft{}Fright, scare\textquotedblright{}.) Say it once more.
\end{tcolorbox}
\section[effrayer]{effrayer — to frighten, to scare}

\label{lesson:FR-C470-effrayer}



\begin{quote}
Before the new one: say the French for to irritate, then the French for fright.
\end{quote}

\subsection*{You'll want to know: effrayer}
\textbf{effrayer} — \textquotedblleft{}to frighten, to scare\textquotedblright{}.

\begin{quote}
L'orage a effrayé le chien toute la nuit.
\end{quote}

\begin{quote}
\emph{The storm frightened the dog all night.}
\end{quote}

\emph{Être effrayé} — to be frightened. \emph{Effrayant} — frightening.

\begin{cousinweb}
\emph{effrayer} comes from Old French \emph{esfreer} \textquotedblleft{}to disturb, to take the peace away\textquotedblright{} (EtymDB), built on \emph{es-} \textquotedblleft{}out\textquotedblright{} and \emph{freer} \textquotedblleft{}to make peace\textquotedblright{}, which is usually traced to Frankish \emph{fridu} \textquotedblleft{}peace\textquotedblright{}.
\end{cousinweb}

\begin{grammarlens}[title={What you've built}]
One more word for telling, explaining and describing at B1.
\end{grammarlens}

\subsection*{Your turn}
\begin{itemize}
\raggedright
  \item \emph{Say it:} \emph{effrayer}
  \item \emph{Say it:} \emph{effrayer}, once more
  \item \emph{Say it:} a sentence of your own with \emph{effrayer}
  \item \emph{Recall:} say the French for to irritate, then the French for fright, then say \emph{effrayer} again
\end{itemize}

\begin{tcolorbox}[breakable,colback=teal!4,colframe=teal!35!black,title={Before you move on}]
What does effrayer mean? (\textquotedblleft{}To frighten, to scare\textquotedblright{}.) Say it once more.
\end{tcolorbox}
\section[sursauter]{sursauter — to jump, to start (with surprise)}

\label{lesson:FR-C470-sursauter}



\begin{quote}
Before the new one: say the French for fright, then the French for to frighten.
\end{quote}

\subsection*{You'll want to know: sursauter}
\textbf{sursauter} — \textquotedblleft{}to jump, to start (with surprise)\textquotedblright{}.

\begin{quote}
J'ai sursauté quand la porte a claqué.
\end{quote}

\begin{quote}
\emph{I jumped when the door slammed.}
\end{quote}

A sudden jump of the body, from surprise or fear.

\begin{cousinweb}
\emph{sursauter} is built from \emph{sur} and \emph{sauter} \textquotedblleft{}to jump\textquotedblright{}, which is inherited from Latin \emph{saltāre} \textquotedblleft{}to dance, to leap\textquotedblright{} (Lewis \& Short).
\end{cousinweb}

\begin{grammarlens}[title={What you've built}]
One more word for telling, explaining and describing at B1.
\end{grammarlens}

\subsection*{Your turn}
\begin{itemize}
\raggedright
  \item \emph{Say it:} \emph{sursauter}
  \item \emph{Say it:} \emph{sursauter}, once more
  \item \emph{Say it:} a sentence of your own with \emph{sursauter}
  \item \emph{Recall:} say the French for fright, then the French for to frighten, then say \emph{sursauter} again
\end{itemize}

\begin{tcolorbox}[breakable,colback=teal!4,colframe=teal!35!black,title={Before you move on}]
What does sursauter mean? (\textquotedblleft{}To jump, to start (with surprise)\textquotedblright{}.) Say it once more.
\end{tcolorbox}
\section[le soulagement]{le soulagement — relief}

\label{lesson:FR-C470-soulagement}



\begin{quote}
Before the new one: say the French for to frighten, then the French for to jump.
\end{quote}

\subsection*{You'll want to know: le soulagement}
\textbf{le soulagement} — \textquotedblleft{}relief\textquotedblright{}.

\begin{quote}
Quel soulagement ! Ma fille est bien arrivée.
\end{quote}

\begin{quote}
\emph{What a relief! My daughter has arrived safely.}
\end{quote}

The verb is \emph{soulager} \textquotedblleft{}to relieve\textquotedblright{}; \emph{je suis soulagé} — I am relieved.

\begin{cousinweb}
\emph{le soulagement} is built on \emph{soulager} \textquotedblleft{}to relieve\textquotedblright{}, which comes from Latin \emph{sublevāre} \textquotedblleft{}to lift up, to lighten\textquotedblright{} (Lewis \& Short), from \emph{levāre} \textquotedblleft{}to raise\textquotedblright{}.
\end{cousinweb}

\begin{grammarlens}[title={What you've built}]
One more word for telling, explaining and describing at B1.
\end{grammarlens}

\subsection*{Your turn}
\begin{itemize}
\raggedright
  \item \emph{Say it:} \emph{le soulagement}
  \item \emph{Say it:} \emph{le soulagement}, once more
  \item \emph{Say it:} a sentence of your own with \emph{le soulagement}
  \item \emph{Recall:} say the French for to frighten, then the French for to jump, then say \emph{le soulagement} again
\end{itemize}

\begin{tcolorbox}[breakable,colback=teal!4,colframe=teal!35!black,title={Before you move on}]
What does le soulagement mean? (\textquotedblleft{}Relief\textquotedblright{}.) Say it once more.
\end{tcolorbox}
\section[rassurer]{rassurer — to reassure}

\label{lesson:FR-C470-rassurer}



\begin{quote}
Before the new one: say the French for to jump, then the French for relief.
\end{quote}

\subsection*{You'll want to know: rassurer}
\textbf{rassurer} — \textquotedblleft{}to reassure\textquotedblright{}.

\begin{quote}
Le médecin a rassuré mes parents.
\end{quote}

\begin{quote}
\emph{The doctor reassured my parents.}
\end{quote}

\emph{Je suis rassuré} — I feel reassured.

\begin{cousinweb}
\emph{rassurer} comes from Old French \emph{rasseürer} (EtymDB), built from \emph{re-} and \emph{asseürer} \textquotedblleft{}to make safe\textquotedblright{}, which goes back to Latin \emph{sēcūrus} \textquotedblleft{}free from care\textquotedblright{} (Lewis \& Short).
\end{cousinweb}

\begin{grammarlens}[title={What you've built}]
One more word for telling, explaining and describing at B1.
\end{grammarlens}

\subsection*{Your turn}
\begin{itemize}
\raggedright
  \item \emph{Say it:} \emph{rassurer}
  \item \emph{Say it:} \emph{rassurer}, once more
  \item \emph{Say it:} a sentence of your own with \emph{rassurer}
  \item \emph{Recall:} say the French for to jump, then the French for relief, then say \emph{rassurer} again
\end{itemize}

\begin{tcolorbox}[breakable,colback=teal!4,colframe=teal!35!black,title={Before you move on}]
What does rassurer mean? (\textquotedblleft{}To reassure\textquotedblright{}.) Say it once more.
\end{tcolorbox}
